\documentclass[11pt]{article}

% ---------------------------------------------------------------------------
%  reharness: A Reasoning Compiler for AI Workflows
%  Technical / scientific documentation of the runtime and compiler (v0.25.0).
%  Build:  pdflatex reharness.tex  (twice, for cross-references)
% ---------------------------------------------------------------------------

\usepackage[a4paper,margin=1in]{geometry}
\usepackage{amsmath,amssymb,amsthm}
\usepackage{mathtools}
\usepackage{enumitem}
\usepackage{booktabs}
\usepackage{xcolor}
\usepackage{listings}
\usepackage{hyperref}
\hypersetup{colorlinks=true,linkcolor=blue!55!black,citecolor=green!45!black,urlcolor=blue!55!black}

% ---- theorem environments -------------------------------------------------
\theoremstyle{plain}
\newtheorem{theorem}{Theorem}[section]
\newtheorem{proposition}[theorem]{Proposition}
\newtheorem{lemma}[theorem]{Lemma}
\newtheorem{corollary}[theorem]{Corollary}
\theoremstyle{definition}
\newtheorem{definition}[theorem]{Definition}
\newtheorem{principle}[theorem]{Principle}
\newtheorem{invariant}[theorem]{Invariant}
\theoremstyle{remark}
\newtheorem{remark}[theorem]{Remark}

% ---- listings -------------------------------------------------------------
\definecolor{codebg}{gray}{0.96}
\lstset{
  basicstyle=\ttfamily\small,
  backgroundcolor=\color{codebg},
  frame=single, rulecolor=\color{gray!40},
  breaklines=true, columns=fullflexible,
  keepspaces=true, showstringspaces=false,
  xleftmargin=4pt, xrightmargin=4pt,
}

% ---- shorthands -----------------------------------------------------------
\newcommand{\reharness}{\textsc{reharness}}
\newcommand{\code}[1]{\texttt{#1}}
\newcommand{\state}[1]{\textsf{#1}}
\newcommand{\transition}{\ensuremath{\delta}}
\newcommand{\Nat}{\mathbb{N}}
\newcommand{\given}{\,\vert\,}
\newcommand{\powerset}{\mathcal{P}}
\newcommand{\meet}{\sqcap}
\newcommand{\joinop}{\sqcup}
% Indexed meet: \bigsqcap is not in base LaTeX; use the standard big operator that
% matches our concrete lattice (intersection = must, union = may).
\newcommand{\bigmeet}{\bigcap}

\title{\reharness: A Reasoning Compiler for AI Workflows\\[4pt]
\large Formal description of the runtime and compiler (v0.25.0)}
\author{reharness project}
\date{\today}

\begin{document}
\maketitle

\begin{abstract}
\reharness{} compiles a natural-language request into a runnable, finite-state
AI workflow. It is a \emph{reasoning compiler}: the nondeterministic intelligence
of large language models (LLMs) is confined to a one-time compilation phase and to
clearly delimited \emph{agent leaves} of an otherwise fully deterministic execution
skeleton. The system has two layers, both finite-state machines (FSMs). The
\emph{runtime} executes a compiled FSM and is a named restriction of a standard,
well-studied model: a \emph{deterministic hierarchical Moore-action transducer with
run-to-completion semantics}. The \emph{compiler} turns a request into a single
declarative artifact---a graph with per-node behavioural contracts---and lowers it
to executable TypeScript through deterministic graph passes, validated by a static
analyzer built on two textbook engines (transitive-closure reachability and the
Kam--Ullman monotone data-flow framework). Inter-stage data flow is \emph{derived}
from graph topology via a polyhedral instance-wise model rather than declared.
We give the formal models, state and prove the determinism and termination
guarantees, and present the compiler's design as a small set of principles from
which the mechanical machinery is derived. We relate the system to prior work on
LLM compilation (DSPy, LLM+P, ``Compiled AI''), statecharts, monotone data-flow
analysis, and the polyhedral model.
\end{abstract}

\tableofcontents

% ===========================================================================
\section{Introduction}
% ===========================================================================

LLM-based agent frameworks typically invoke a model at \emph{every} step of a
workflow, conflating two distinct concerns: \emph{deciding what to do} (which
benefits from reasoning) and \emph{executing a well-specified procedure} (which
does not). The conflation is the source of three well-documented problems:
nondeterminism (outputs vary run to run, even at temperature $0$), cost (tokens
scale with execution volume), and unauditability (no fixed artifact to inspect).

\reharness{} takes the opposite stance, shared with a small body of recent work
on treating ``LLM as compiler, not interpreter'': intelligence is spent
\emph{once}, at compile time, to produce a deterministic artifact that is then
executed by conventional software. Where intelligence is genuinely irreducible at
run time---a judgement over unstructured input---it is isolated into an
\emph{agent leaf}: a single LLM call that is one state of a deterministic FSM.

\paragraph{Two layers.} The architecture is two finite-state machines:
\begin{itemize}[nosep]
  \item the \textbf{runtime} (\S\ref{sec:runtime}, \code{src/runtime/}): executes a
        compiled FSM---states are agent/code/$\dots$ actions, transitions are typed
        and guarded;
  \item the \textbf{compiler} (\S\ref{sec:compiler}, \code{src/compiler/}): the
        \code{/generate} meta-pipeline turns a request into a \code{skeleton.xml}
        and lowers it to a runnable TypeScript pipeline; the analyzer in
        \code{src/compiler/analysis/} statically validates it.
\end{itemize}
The compiler is itself a \reharness{} pipeline executed by the same runtime; the
system is self-hosting.

\paragraph{Design ethos.} Both layers are \emph{named restrictions of standard
models}, not ad hoc constructions. The runtime is a restriction of the
Harel/UML hierarchical state machine; the analyzer is a family of instances of two
standard engines; the data-flow model is the polyhedral iteration space. The
guiding engineering principle is \emph{simplify until it breaks determinism}: the
codebase has repeatedly \emph{removed} machinery (declared data flow, review loops,
cross-cutting spec blocks) when a seam cost more than the check it bought. This
document records the resulting theory.

% ===========================================================================
\section{The Runtime}\label{sec:runtime}
% ===========================================================================

The runtime (\code{src/runtime/fsm.ts}) executes a compiled pipeline. Formally it
is a \textbf{deterministic hierarchical Moore-action transducer driven by
run-to-completion (RTC) steps}: a deterministic finite control with a recursion
stack for composite states (a restricted pushdown structure), where each state's
action is a self-contained computation that emits exactly one event. It is a
deliberate restriction of the UML/Harel hierarchical state machine (HSM)---the same
control topology, minus the machinery whose nondeterminism we reject
(\S\ref{sec:restrictions}).

\subsection{Syntax of a pipeline}

\begin{definition}[Pipeline]\label{def:pipeline}
A pipeline is a tuple $P = (Q, q_0, F, \tau, \delta_{\mathrm{on}}, G, \kappa)$ where
\begin{itemize}[nosep]
  \item $Q$ is a finite set of \emph{states} and $q_0 \in Q$ the \emph{initial} state;
  \item $F \subseteq Q$ are \emph{final} states, each labelled \code{success} or \code{error};
  \item $\tau : Q \to \{\state{agent}, \state{code}, \state{interactive}, \state{switch},
        \state{set}, \state{parallel}, \state{loop}, \state{call}, \state{wait}, \state{final}\}$
        assigns each state a \emph{type};
  \item $\delta_{\mathrm{on}} : Q \times \Sigma \rightharpoonup (Q \times G)^{*}$ maps a state and an
        \emph{event} $e \in \Sigma$ to an ordered list of guarded targets;
  \item $G$ is the set of \emph{guards} (boolean expressions over the data store, \S\ref{sec:guards});
  \item $\kappa$ assigns composite states their structural parameters (branch/step bodies, join target,
        iteration bound, etc.).
\end{itemize}
\end{definition}

States divide into \emph{active} states (\state{agent}, \state{code},
\state{interactive}), whose semantics is an \emph{action}; \emph{control} states
(\state{switch}, \state{set}); \emph{composite} states (\state{parallel},
\state{loop}, \state{call}); the external-signal state \state{wait}; and
\state{final} states.

\subsection{Moore-action semantics and the run-to-completion step}

\begin{definition}[Moore action]
Each active state $q$ carries an \emph{entry action} $a_q$ that runs to completion
and \emph{emits exactly one event} $e \in \Sigma$ equal to the \emph{outcome of its
own computation}. The emitted event is not a function of an input symbol read
mid-step (that would be a Mealy machine); it is the result of $a_q$.
\end{definition}

A configuration is a pair $(q, \sigma)$ where $q$ is the active state and $\sigma$
is the \emph{data store} (an in-memory key/value map \code{ctx.data}, plus the
read-only \code{config} and the per-stage workspace; \S\ref{sec:dataflow}). The
runtime executes one RTC step at a time.

\begin{definition}[RTC step]\label{def:rtc}
A run-to-completion step from configuration $(q,\sigma)$ with $\tau(q)$ active:
\begin{enumerate}[nosep]
  \item run $a_q$ to completion, producing event $e$ and an updated store $\sigma'$;
  \item compute the next state $q' = \transition(q, e, \sigma')$ (\S\ref{sec:delta});
  \item the next step begins only after the current step has fully settled.
\end{enumerate}
There is no event queue, no preemption, and no interleaving: exactly one event
exists per step, so the RTC requirement ``no internal concurrency within a step''
holds trivially. \state{wait} states are the only source of externally-timed events.
\end{definition}

\subsection{The transition function}\label{sec:delta}

\begin{definition}[Transition function]\label{def:delta}
Given a state $q$, an emitted event $e$, and a store $\sigma$, let
$\delta_{\mathrm{on}}(q,e) = \langle (t_1,g_1), \dots, (t_k,g_k)\rangle$ be the
ordered guarded-target list. Then
\[
  \transition(q,e,\sigma) \;=\; t_{j}, \qquad
  j = \min\{\, i \given g_i(\sigma) = \text{true} \,\},
\]
i.e.\ the target of the \emph{first} guard that holds in document order. If
$\delta_{\mathrm{on}}(q,e)$ is undefined, or no $g_i$ holds, the runtime
\emph{fails loud}: it halts with a diagnostic rather than stalling silently.
\end{definition}

\begin{definition}[Totality by completion]\label{def:total}
A pipeline is \emph{total} if for every reachable configuration $(q,\sigma)$ and
every event $e$ that $a_q$ can emit, $\transition(q,e,\sigma)$ is defined and some
guard holds. Codegen and lint enforce structural conditions sufficient for
totality (every emitted event has a transition; \state{switch} states have a
default arm; an \state{error} final exists for the synthesized \code{ERROR} edge);
where they cannot, $\transition$ fails loud (Def.~\ref{def:delta}), so the
runtime never silently stalls. We treat ``fail loud'' itself as making $\transition$
total into $Q \cup \{\bot\}$ with $\bot$ an observable error halt.
\end{definition}

\subsection{The guard language}\label{sec:guards}

Guards $g \in G$ are drawn from a deliberately tiny, side-effect-free expression
sublanguage, so that they are statically analyzable and cannot themselves introduce
nondeterminism or arbitrary computation. The grammar admits: the root identifiers
\code{config}, \code{data}, \code{retries} with member access; the comparison and
boolean operators ${==}\ {!=}\ {<}\ {\le}\ {>}\ {\ge}\ \&\&\ \vert\vert\ !$ and the
arithmetic operators ${+}\ {-}\ {*}\ {/}\ \%$; string/number/boolean/\code{null}
literals; and array literals. It \emph{rejects} function calls, assignment, the
ternary operator, increment/decrement, bitwise operators, and regular expressions.
A retry guard has the special form $\code{retries:}k < N$ and reads a bounded retry
counter $k$. The compiler lowers a guard to a pure TypeScript arrow-function body
(root identifiers prefixed with the context object); the runtime never evaluates a
guard as a string. Rejecting calls is also a safety property: guard text originates
in LLM-authored XML and is embedded verbatim into generated code, so an admitted
call would be a code-execution surface (\S\ref{sec:limitations}).

\subsection{Hierarchy: composites as RTC sub-computations}

Composite states run as nested RTC sub-computations with their own recursion frame
and return a single completion to the parent; the parent transition treats the
whole composite as one Moore action.

\paragraph{\state{parallel} (fork--join).}
A \state{parallel} state evaluates an expression \code{over} to an array
$[x_0,\dots,x_{m-1}]$, runs a \emph{branch} body once per item (up to a concurrency
cap), and after all branches settle transitions to its \code{join}. Agent branches
run with \emph{real OS-process parallelism} (each spawns a separate \code{pi}
process); pure-code branches run cooperatively on the single event loop. The join
reads the per-branch results $\code{data.branches} = [\{i, x_i, \mathit{dir}_i,
\mathit{ok}_i\}]$ as a barrier.

\begin{invariant}[No shared writes in a branch]\label{inv:branch}
A \state{parallel} branch must not write the shared store \code{ctx.data}.
Branches communicate only through their isolated output directories; cross-branch
data is read by the join after the barrier.
\end{invariant}

\paragraph{\state{loop} (bounded iteration).}
A \state{loop} runs an ordered list of \emph{step} states once per iteration. It
carries a mandatory hard bound \code{max} $\in \Nat$ and an optional early-exit
predicate \code{exit}. The runtime maintains two distinct iteration-space
quantities (\S\ref{sec:loopcounters}).

\paragraph{\state{call} (sub-machine).}
A \state{call} invokes another compiled pipeline as a sub-pipeline and transitions
by its terminal status (\code{success}/\code{error}).

\subsection{Determinism}

\begin{theorem}[Determinism]\label{thm:determinism}
Fix a pipeline $P$ and the outputs of all agent leaves. Then the observable result
of executing $P$---the sequence of states visited, the final store, and the set of
workspace files written---is a deterministic function of the initial store, up to
the commutative merge of \state{parallel} branch outputs.
\end{theorem}

\begin{proof}[Proof sketch]
We show every source of choice is resolved deterministically.
(i) Within a step, $a_q$ is run to completion and emits a single event
(Def.~\ref{def:rtc}); the only nondeterminism inside $a_q$ is the agent leaf, which
is fixed by hypothesis. (ii) $\transition$ is a function, not a relation: at most
one transition fires per $(q,e)$ because guards are scanned in document order and
the \emph{first} true guard wins (Def.~\ref{def:delta}); there is no
``highest-priority'' ambiguity and no simultaneously-enabled race. (iii) Composites
return a single completion, so the parent sees one Moore action. (iv) For
\state{parallel}: branches cannot write shared \code{ctx.data}
(Inv.~\ref{inv:branch}) and the join is a barrier, so the only order-dependence is
the order in which branch outputs are merged; since each branch writes a disjoint,
index-addressed output directory (\S\ref{sec:dataflow}), the merge is commutative
and the observable result is independent of completion order. Composing
(i)--(iv) over the RTC step sequence yields a deterministic transducer up to that
commutative merge.
\end{proof}

\begin{remark}
Theorem~\ref{thm:determinism} is the design's central contract: the FSM is the
\emph{deterministic skeleton} around nondeterministic agent leaves. ``Compiled AI''
in the sense of Trooskens et al.\ pushes the leaves to zero (no run-time LLM); a
\reharness{} pipeline keeps leaves where reasoning is irreducible but keeps the
\emph{harness} around them deterministic. The fully-deterministic case
(zero agent states) is a reachable point of the design, not a separate mode
(\S\ref{sec:amortization}).
\end{remark}

\subsection{Termination}

\begin{invariant}[Loops are bounded]\label{inv:max}
Every \state{loop} state declares a hard bound $\code{max} \in \Nat$, $\code{max}
\ge 1$ (enforced by lint, \S\ref{sec:lint}). The early-exit predicate \code{exit}
is an optional refinement, never a substitute for the bound.
\end{invariant}

\begin{theorem}[Termination]\label{thm:termination}
Every execution of a pipeline $P$ in which each agent leaf terminates halts in a
finite number of RTC steps.
\end{theorem}

\begin{proof}[Proof sketch]
Define a well-founded rank on configurations. A \state{loop} contributes a counter
bounded above by its \code{max} (Inv.~\ref{inv:max}); it strictly increases each
iteration and forces exit at \code{max}, so a loop performs at most $\code{max}$
iterations. A \state{parallel} performs $m = |\code{over}|$ branch executions, $m$
finite. A \state{call} recurses into a strictly smaller sub-pipeline (no cyclic
call graph by construction). Top-level cycles in the control graph are admissible
only when gated by a \code{retries} counter, which is bounded by its
\code{retries-max} (a retry edge is the only way to re-enter an earlier state, and
the counter is monotone). Thus every back-edge is bounded by a static counter, and
the lexicographic product of these counters is a well-founded measure that strictly
decreases along every step that could otherwise recur. Finitely many non-recurring
steps plus a well-founded measure on recurring ones gives finiteness, provided each
$a_q$ (in particular each agent leaf) terminates.
\end{proof}

\subsection{Loop counters: coordinate versus cardinality}\label{sec:loopcounters}

A subtle but instructive point. A loop exposes two \emph{categorically distinct}
quantities to the data store, and conflating them is a latent off-by-one bug.

\begin{definition}[Loop coordinate and cardinality]
For a loop with iterations indexed $0,1,\dots$:
\begin{itemize}[nosep]
  \item \code{data.iteration} is the \textbf{coordinate}: the $0$-based index of the
        \emph{current} iteration. It is a component of the instance vector
        (\S\ref{sec:dataflow}), valid \emph{inside} a step. The early-exit
        predicate \code{exit} sees it as the index of the iteration that just ran.
  \item \code{data.iterations} is the \textbf{cardinality}: the number of iterations
        that ran. It is the loop's return value, available \emph{after} the loop.
\end{itemize}
\end{definition}

The distinction mirrors the difference between an array index $i$ and the array
length $|A|$. ``How many rounds ran'' is the cardinality \code{data.iterations}; it
must never be computed as $\code{data.iteration}+1$, and \code{data.iteration} must
not be read after the loop. The runtime restores the coordinate (not the
cardinality) into the shared scalar after each iteration so that a nested loop step
that overwrote it does not corrupt the enclosing loop's view.

\subsection{Restrictions relative to the full UML HSM}\label{sec:restrictions}

\reharness{}'s runtime deliberately omits, from the full statechart model:
\begin{itemize}[nosep]
  \item \textbf{orthogonal regions with event broadcast}---replaced by
        \state{parallel} fork--join with \emph{no} inter-branch signalling.
        Broadcast is the principal source of statechart nondeterminism and
        state-space blow-up;
  \item \textbf{Mealy transition actions and entry/exit action chains}---all
        behaviour lives in Moore-state actions; transitions are pure control;
  \item \textbf{internal/local transitions, deep history, completion-event
        races}---omitted.
\end{itemize}
What remains is exactly the analyzable core: a control graph, a recursion stack,
and per-state actions emitting one event. This is what makes
Theorems~\ref{thm:determinism} and \ref{thm:termination} hold by construction.

% ===========================================================================
\section{The Data-Flow and Workspace Model}\label{sec:dataflow}
% ===========================================================================

A central design decision: \emph{inter-stage data flow is derived from graph
topology, never declared}. There is no \code{produces}/\code{consumes} attribute,
no path syntax, no hand-built wiring. We motivate this historically, then give the
polyhedral model that realizes it.

\subsection{Why derived, not declared}

Two earlier designs declared data flow---first file artifacts with
\code{artifact=}/\code{requires=} paths, then a \code{produces}/\code{consumes}
layer---and both \emph{relocated} the drift seam rather than removing it: path
strings drift from reality, and declaration-vs-prose drift recurs at the
fan-in of parallel stages. The realization:

\begin{principle}[An edge is the data contract]\label{prin:edge}
A graph edge $A \to B$ already \emph{is} the statement that $B$ may read $A$'s
output. Therefore the author declares nothing about wiring; the compiler derives it
from edges.
\end{principle}

\subsection{Three channels}

Data moves through three channels:
\begin{enumerate}[nosep]
  \item \textbf{\code{config}}---the external CLI interface (\S\ref{sec:inputs}),
        read-only, \emph{declared} (it has no internal producer to derive from);
  \item \textbf{\code{ctx.data}}---in-memory scalars, written only by \state{code}
        and \state{set} states, read by guards / \code{over} / \code{exit} /
        \code{model-expr};
  \item \textbf{per-stage workspace directories}---files; every stage gets its own
        output directory and reads its visible producers' directories.
\end{enumerate}

\subsection{The polyhedral instance-wise model}

The workspace model is the \emph{polyhedral iteration space} (Feautrier's
instance-wise array data-flow analysis): a stage executes once per point in its
iteration space, and a read sees the producer instances on its shared-scope prefix,
collected over the axes it has exited.

\begin{definition}[Enclosing scope / iteration vector]\label{def:scope}
For a state $n$, its \emph{enclosing scope} $\mathrm{scope}(n) = \langle c_1,
\dots, c_d \rangle$ is the chain of composite states enclosing it, outermost first
(a \state{parallel} contributes a branch axis, a \state{loop} an iteration axis).
A single \emph{instance} of $n$ is addressed by an \emph{instance vector}
$\vec{\imath} = (i_1, \dots, i_d) \in \Nat^d$, one index per axis. Its output
directory is
\[
  \mathrm{dir}(n, \vec{\imath}) \;=\; \langle\mathit{runDir}\rangle/\code{work}/n/i_1/i_2/\cdots/i_d .
\]
A top-level stage has $d = 0$ and a single directory.
\end{definition}

\begin{definition}[Visible producers]\label{def:visible}
The producers visible to $n$ are its \emph{ancestor producers}: states of type
\state{agent}/\state{code}/\state{interactive} backward-reachable from $n$ through
the role-aware successor relation (\S\ref{sec:engines}). Each producer $p$ is
classified by output cardinality relative to $n$. Let $\ell =
|\mathrm{lcp}(\mathrm{scope}(n),\mathrm{scope}(p))|$ be the length of the longest
common prefix of the two scopes. Then
\[
  p \in
  \begin{cases}
    \mathit{single}(n) & \text{if } |\mathrm{scope}(p)| = \ell, \\[2pt]
    \mathit{list}(n)   & \text{otherwise.}
  \end{cases}
\]
\end{definition}

Intuitively: $p$ is \emph{single} iff every composite enclosing $p$ also encloses
$n$ (so $p$'s instance is fixed by $n$'s surrounding context---linear chains,
same-scope reads, loop-carried reads). $p$ is a \emph{list} iff $n$ has exited at
least one composite enclosing $p$, in which case $n$ sees a \emph{collection} of
$p$'s instances over the exited axes: one directory per branch (an exited
\state{parallel} is a \emph{map}) and/or per iteration (an exited \state{loop} is a
\emph{scan} / history).

\begin{proposition}[One law for map, scan, and nesting]
Definition~\ref{def:visible} yields correct visibility and cardinality uniformly
for: linear chains (including grand-ancestor skips), \state{parallel} fan-out/in
(map), \state{loop} cross-iteration history (scan) and loop-carried reads, and
nested composites at arbitrary depth (the instance vector grows one axis per
composite).
\end{proposition}

At run time, three accessors expose this to a stage:
\code{c.out()} = its own instance directory; \code{c.dir(p)} = a single
shared-scope producer instance (loop-carried-aware: the latest existing instance
$\le$ the current index); \code{c.dirs(p)} = the collection over exited axes.
Producer-write, the recorded \code{data.branches[i].dir}, and consumer-read all
derive from $(\,\text{stage}, \vec{\imath}\,)$, so they cannot drift.

\begin{remark}[Boundary]
Topology derives \emph{wiring} and \emph{cardinality} (which instances are
visible), not \emph{shape} (the schema of a file---a behavioural property checked by
the \state{polish} pass) nor \emph{selectivity} (which subset a stage should read,
governed by the contract prose). Conditional producers on alternate \state{switch}
arms are offered permissively (no ``produced on every path'' guarantee for files;
only \code{ctx.data} carries the must-write check of \S\ref{sec:defassign}).
\end{remark}

% ===========================================================================
\section{The Static Analyzer}\label{sec:compiler-analysis}
% ===========================================================================

Every static check \reharness{} performs is an \emph{instance} of one of two
reusable engines (\code{src/compiler/analysis/framework.ts}). New checks are added
as instances; the codebase does not hand-roll bespoke fixpoints.

\subsection{The two engines}\label{sec:engines}

\begin{definition}[Reachability engine]
$\mathrm{reachableFrom}(S, \mathrm{next})$ computes the transitive closure of the
set $S$ under the successor relation $\mathrm{next}$ by worklist BFS. The lattice is
trivial (a reachable set).
\end{definition}

\begin{definition}[Monotone data-flow engine]\label{def:kamullman}
$\mathrm{solveMonotoneSets}$ is the \textbf{Kam--Ullman monotone data-flow
framework} specialised to the lattice $L = (\powerset(K), \subseteq)$ of subsets of
a finite key set $K$, with a gen-only forward transfer function
\[
  \mathrm{OUT}[n] = \mathrm{gen}(n) \cup \mathrm{IN}[n], \qquad
  \mathrm{IN}[n] = \!\!\bigmeet_{p \in \mathrm{pred}(n)}\!\! \mathrm{OUT}[p],
\]
iterated to a fixpoint, where $\meet$ is set intersection
(\textbf{must} analysis) or set union (\textbf{may} analysis).
\end{definition}

\begin{proposition}[Convergence]\label{prop:converge}
$\mathrm{solveMonotoneSets}$ terminates at the least fixpoint. The transfer
functions are monotone and the lattice $\powerset(K)$ has finite height $|K|$;
hence by the Kam--Ullman theorem the chaotic iteration converges.
\end{proposition}

\subsection{The checks, as instances}

\begin{table}[h]
\centering
\small
\begin{tabular}{@{}lll@{}}
\toprule
\textbf{Check} & \textbf{Engine instance} & \textbf{Module} \\
\midrule
reachable-from-initial   & forward reachability from $q_0$           & \code{semantic.ts} \\
can-reach-a-final        & backward reachability from $F$ (co-reach.) & \code{semantic.ts} \\
definite assignment      & forward \emph{must} (meet $=$ intersect)  & \code{dataflow.ts} \\
\code{visibleProducers}  & backward may-reach.\ $+$ cardinality      & \code{graph.ts} \\
config-flow              & read-set $\subseteq$ declared inputs      & \code{dataflow.ts} \\
grammar / well-formedness& the format's ``type system''             & \code{lint.ts} \\
contract coverage        & every active state has a \code{<contract>} & \code{semantic.ts} \\
\bottomrule
\end{tabular}
\caption{Each static check is a named instance of an engine in \S\ref{sec:engines}.}
\label{tab:checks}
\end{table}

\subsection{Reachability and dead-ends}

Reachable-from-initial is $\mathrm{reachableFrom}(\{q_0\}, \mathrm{succ})$; any
state not in the result is unreachable. Can-reach-a-final is
$\mathrm{reachableFrom}(F, \mathrm{succ}^{-1})$ over reverse edges (co-reachability);
any state not in the result is a dead end. The successor relation is
\emph{role-aware}: a \state{parallel} branch or \state{loop} step routes through its
parent composite's join, not its own \code{<on>} edges.

\subsection{Definite assignment over \code{ctx.data}}\label{sec:defassign}

\begin{definition}[Definite assignment]
A read of \code{data.k} at state $n$ is \emph{safe} if \code{k} is written on
\emph{every} path reaching $n$. This is the textbook forward-\emph{must} data-flow
problem: the instance of Def.~\ref{def:kamullman} with $\meet = {}$intersection and
$\mathrm{gen}(n)$ the keys $n$ definitely writes. Let $\mathrm{IN}[n]$ be the
fixpoint. A read of \code{data.k} at $n$ is flagged unless $\code{k} \in
\mathrm{IN}[n]$.
\end{definition}

\begin{proposition}[Soundness of the must-check]
If the check does not flag a read of \code{data.k} at $n$, then \code{k} is written
on every control path from $q_0$ to $n$; hence \code{data.k} is defined whenever $n$
executes. Equivalently, the check flags a read only when a writer-free path provably
exists.
\end{proposition}

A refinement distinguishes \emph{entry-reads} (declared reads; the \code{over}
expression of a \state{parallel}, evaluated at entry; the value expressions of a
\state{set}) from \emph{after-reads} (the guards on the state's own outgoing
transitions, evaluated \emph{after} the entry action has run). Entry-reads are
checked against $\mathrm{IN}[n]$; after-reads against $\mathrm{IN}[n] \cup
\mathrm{gen}(n)$, since a node's own writes are available to its own outgoing guards.
This eliminates a false positive on the common pattern ``a \state{code} state writes
\code{data.x} then routes on \code{expr:data.x}''.

% ===========================================================================
\section{The External Interface}\label{sec:inputs}
% ===========================================================================

The one thing about data that \emph{is} declared is the CLI interface, because it
is external---it comes from the user, not from any graph edge, so there is no
producer to derive it from.

\begin{principle}[Declare the external, derive the internal]\label{prin:law}
Internal data flow is derived from graph edges (Principle~\ref{prin:edge}). The
external interface (CLI inputs) is declared, in a per-pipeline \code{<inputs>}
block. This is the irreducible signature of the pipeline, not a return to declared
data flow.
\end{principle}

Each \code{<arg>} declares a \code{config.X} field; codegen generates the argument
parser from the declaration (positional vs.\ flag, type coercion among
\code{string}/\code{list}/\code{number}/\code{bool}, defaults, required-ness). Two
fields are always provided and never declared: \code{config.target} (the working
directory) and \code{config.input} (positional arguments joined).

\begin{definition}[Config-flow check]
Let $\mathrm{Decl}$ be the declared inputs $\cup\, \{\code{target}, \code{input}\}$
and $\mathrm{Used}$ the set of \code{config.X} fields read anywhere (skeleton
expressions: guards, \code{over}, \code{exit}, \code{model-expr}, \code{set} values;
and the generated code: \code{c.config.X}). The check requires $\mathrm{Used}
\subseteq \mathrm{Decl}$.
\end{definition}

\begin{corollary}[CLI correctness by construction]
If a pipeline passes the config-flow check, every \code{config.X} the pipeline reads
is parsed from the command line by the generated parser. Hence ``compiles''
implies ``CLI wiring is correct''.
\end{corollary}

% ===========================================================================
\section{The Compiler Pipeline}\label{sec:compiler}
% ===========================================================================

The compiler (\code{/generate}) is a self-hosted \reharness{} pipeline: it is an
FSM executed by the runtime of \S\ref{sec:runtime}. Its stages are:
\begin{center}\small
\state{maybe\_research} $\to$ \state{research} $\to$ \state{prd} $\to$
\state{review\_prd} \emph{(approval)} $\to$ \state{design} $\to$ \state{construct}
$\to$ \state{fill\_prompts} $\to$ \state{check\_dataflow} $\to$ \state{polish}
$\to$ \state{verify} $\to$ \state{done}.
\end{center}

\subsection{The lean principle and the single artifact}

\begin{principle}[The LLM authors the minimal artifact]\label{prin:lean}
The LLM writes exactly one artifact---the \code{skeleton.xml}: the \emph{graph}
(states $+$ typed transitions) plus a per-node behavioural \code{<contract>}
(CDATA). Everything mechanical (the argument parser, the per-stage wiring, the
command and library scaffolding) is a deterministic graph pass. There is no
\code{produces}/\code{consumes}, no \code{<spec>}, no hand-built path.
\end{principle}

\subsection{The human approves intent, never the graph}

\begin{principle}[PRD-first]\label{prin:prd}
The single human checkpoint is \state{review\_prd}: the human edits and approves the
\emph{PRD} (the product requirements document: goal, behaviour, acceptance
criteria, scope)---the statement of \emph{intent}. The human never approves the FSM
graph, which is a mechanical realization the compiler owns. Intent is the thing a
human can judge; the realization is the thing the analyzer can check.
\end{principle}

The PRD is also the point of \emph{unification} of input sources: \code{/generate}
maps $\text{prompt} + \text{research} \to \text{PRD}$, and the planned
\code{compile} front-end maps $\text{session-logs} + \text{research} \to
\text{PRD}$; everything downstream of the PRD is reused unchanged.

\subsection{Lowering and verification}

\state{construct} parses and statically validates the skeleton (lint $+$ semantic,
Table~\ref{tab:checks}) and emits the command, the code-state library stubs, and the
agent prompt stubs. \state{fill\_prompts} turns each \code{<contract>} into either
an agent prompt (for \state{agent} states) or a code implementation (for \state{code}
states), in parallel. \state{check\_dataflow} runs the data-flow analyzer
(definite-assignment $+$ config-flow). \state{verify} compiles the generated
TypeScript (\code{tsc --noEmit}) and loads each command's FSM graph. Note that
verification is \emph{static}: the compiler never executes the generated pipeline;
behavioural validation against runs is the province of a separate, future
feedback loop (\S\ref{sec:limitations}).

\subsection{Correction: one corrector in a hot context}\label{sec:lint}

\begin{principle}[Polish, not a review loop]\label{prin:polish}
Correction is a single \state{polish} agent that reviews the whole generated
pipeline against the PRD and fixes the \emph{leaves} (agent prompts, code
implementations) in one hot context---replacing the older
review$\to$fix$\to$re-review loop. \state{redesign} is a rare last resort for
topology changes (when \state{polish} cannot fix a problem in the leaves and
escalates); it stays faithful to the approved PRD---it changes \emph{how}, never
\emph{what}.
\end{principle}

\begin{principle}[Decisions belong to states or agents]\label{prin:decisions}
Routing decisions belong to FSM states (\state{switch}, guarded transitions);
judgement decisions belong to an agent. Neither belongs to imperative glue between
stages. The runtime does not parse, classify, or gate between agent turns.
\end{principle}

% ===========================================================================
\section{Amortization of Reasoning}\label{sec:amortization}
% ===========================================================================

A \state{code} state is \emph{frozen reasoning}: authored once at compile time, then
executed at zero token cost forever. Amortization is the compiler's policy of
realizing a stage as \state{code} rather than \state{agent} whenever the reasoning
is not actually irreducible. Because invariant ``a node is wholly \state{code} or
wholly \state{agent}'' holds (an embedded LLM call inside a code state is rejected
by lint), amortization is a \emph{total} demotion $\state{agent} \to \state{code}$,
and its boundary is objective.

\begin{definition}[Demotability]\label{def:demote}
A stage $n$ is \emph{demotable} to \state{code} iff both:
\begin{enumerate}[label=(\alph*),nosep]
  \item \textbf{structured input}: every producer $n$ reads (its ancestors via
        Def.~\ref{def:visible}) is a \state{code}/\state{set} state, so the bytes
        carry a compiler-authored schema. \emph{This condition is statically
        decidable} from the graph.
  \item \textbf{mechanical contract}: the contract is a total computable function
        over that structure (reformat, filter, sort, dedup, tally, arithmetic,
        extract a field, validate a schema). It is \emph{not} a judgement
        (summarize, rate, classify ambiguous text, cluster ``same issue?'',
        assign severity). \emph{This condition is not statically decidable}---by
        Rice's theorem, the semantic property ``this specification is a pure
        function over the structure'' is undecidable---so it is a judgement, which
        by Principle~\ref{prin:decisions} belongs to the design agent, not to a
        static rule.
\end{enumerate}
\end{definition}

\begin{remark}[Asymmetry and the over-demotion hazard]
Condition (a) is necessary but not sufficient: one can perform judgement over JSON.
Demoting when (b) is false is the worst error---the empirically documented failure
mode in which a regex over an invoice extracts the vendor name at $\sim$20\% accuracy
where an LLM reaches $\sim$80\%: the input looked parseable but the task was
semantic. The policy is therefore conservative: when (b) is in doubt, keep the stage
an \state{agent}. The rule is realized inline in the \state{design} prompt (it costs
no extra pass and no extra tokens); the judgement of (b) is made once, at authoring.
\end{remark}

The two limiting cases are both reachable: a purely generative linear pipeline
(e.g.\ research $\to$ outline $\to$ render) admits no demotion beyond final I/O,
because the chain is agent-fed throughout; a purely mechanical pipeline (e.g.\
validate-then-concatenate two files) demotes to \emph{zero} agent states---a fully
deterministic, zero-token workflow, which is exactly the ``Compiled AI'' limit
reached from the agent-first side.

% ===========================================================================
\section{Relation to Prior Work}
% ===========================================================================

\paragraph{LLM-as-compiler.} DSPy compiles declarative LLM pipelines into optimized
prompts/weights; LLM+P translates natural language into a formal planning
specification (PDDL) delegated to a classical solver; ``Compiled AI'' (Trooskens et
al., 2026) generates deterministic code at a compilation phase and executes it with
zero run-time model calls, with a four-stage generation-and-validation pipeline and
a template/module/prompt-block component library. \reharness{} shares the
compile-once stance but differs in three ways: (i) it compiles \emph{orchestration}
(an analyzable FSM with derived data flow) rather than emitting business code into
fixed templates; (ii) it keeps irreducible reasoning as analyzable agent leaves
rather than forbidding all run-time inference (its ``Code Factory''/bounded-agentic
escape hatch is, in \reharness{}, the primary and only mode, governed by
Def.~\ref{def:demote}); and (iii) the human approves \emph{intent} (the PRD), not a
configuration.

\paragraph{State machines.} The runtime is a named restriction of Harel
\emph{Statecharts} and \emph{UML} run-to-completion semantics, dropping orthogonal
regions and Mealy edge-actions (\S\ref{sec:restrictions}). The \state{parallel}
construct is a fork--join task model, not orthogonal-region concurrency.

\paragraph{Program analysis.} The analyzer is the Kam--Ullman monotone data-flow
framework (reaching definitions, live variables, and definite assignment are its
classical instances) plus transitive-closure reachability; the data-flow/workspace
model is the polyhedral iteration space (Feautrier's instance-wise array data-flow
analysis).

% ===========================================================================
\section{Limitations and Boundaries}\label{sec:limitations}
% ===========================================================================

\begin{itemize}[nosep]
  \item \textbf{Compilation is static by design.} The compiler never runs the
        generated pipeline; \state{verify} checks types and graph loading, not
        behaviour. Behavioural correctness (e.g.\ a consensus rule that is
        well-formed but vacuous, an edge-case in a code state) is observable only
        from runs and is the province of a separate feedback loop (\emph{evolve}),
        connected to the compiler by an external $\text{compile} \to \text{run} \to
        \text{evolve}$ cycle, not by execution inside \code{/generate}.
  \item \textbf{The intent source is variable.} A short prompt leans heavily on the
        (nondeterministic) web-research stage, so the PRD---and hence the
        topology---varies run to run. This is intent variability, not realization
        drift; the human-approves-PRD checkpoint (Principle~\ref{prin:prd}) is the
        anchor against it.
  \item \textbf{Conditional file producers} on alternate \state{switch} arms have no
        must-produce guarantee (only \code{ctx.data} has the must-write check);
        \state{call} sub-pipeline outputs are not modelled as producers.
  \item \textbf{Shape and selectivity are out of scope for topology.} Topology
        derives wiring and cardinality, not file schemas or which subset to read;
        those are behavioural, checked by \state{polish} or stated in contracts.
\end{itemize}

% ===========================================================================
\section{Conclusion}
% ===========================================================================

\reharness{} is a reasoning compiler whose two layers are named restrictions of
standard models: a deterministic hierarchical Moore-action transducer for execution
(deterministic and terminating by construction, Theorems~\ref{thm:determinism} and
\ref{thm:termination}), and a lean, PRD-first compiler whose every static check is
an instance of the Kam--Ullman framework or transitive-closure reachability, and
whose inter-stage data flow is derived from the polyhedral iteration space rather
than declared. The unifying law---\emph{derive the internal, declare the
external}---and the amortization rule---\emph{demote to code iff input is structured
and the contract is mechanical}---let the same pipeline span the full spectrum from
purely generative (all agents) to purely deterministic (zero agents, the Compiled-AI
limit). The design's discipline is to keep that skeleton deterministic and
analyzable, and to spend nondeterministic intelligence only where it is irreducible.

\appendix
\section{Notation}
\begin{center}\small
\begin{tabular}{@{}ll@{}}
\toprule
Symbol & Meaning \\
\midrule
$Q, q_0, F$ & states, initial state, final states \\
$\Sigma$ & event alphabet (\code{DONE}, \code{ERROR}, \code{PASS}, $\dots$) \\
$\transition$ & transition function (Def.~\ref{def:delta}) \\
$\sigma$ & data store (\code{ctx.data} $+$ \code{config} $+$ workspace) \\
$\mathrm{scope}(n)$ & enclosing-composite chain of $n$ (Def.~\ref{def:scope}) \\
$\vec{\imath}$ & instance vector, one index per axis \\
$\mathit{single}(n), \mathit{list}(n)$ & visible producers by cardinality (Def.~\ref{def:visible}) \\
$L = (\powerset(K), \subseteq)$ & data-flow lattice (Def.~\ref{def:kamullman}) \\
$\meet$ & meet ($\cap$ for must, $\cup$ for may) \\
\bottomrule
\end{tabular}
\end{center}

\end{document}
